\documentclass[10pt,a4paper]{article}

\usepackage[margin=1.8cm]{geometry}
\usepackage[T1]{fontenc}
\usepackage[utf8]{inputenc}
\usepackage[portuguese]{babel}
\usepackage{lmodern}
\usepackage{microtype}
\usepackage{array}
\usepackage{amssymb}
\usepackage{booktabs}
\usepackage{xcolor}

\setlength{\parindent}{0pt}
\setlength{\parskip}{0.35em}
\newcommand{\validar}{\textcolor{blue!70!black}{\textbf{A validar}}}

\title{Escalas dos Técnicos de Radiologia\\
Resumo das Restrições para Validação}
\author{Projeto de apoio à decisão para elaboração de escalas}
\date{Versão de \today}

\begin{document}
\maketitle

Este resumo apresenta as regras atualmente implementadas no protótipo. Não
constitui validação jurídica ou institucional. Pretende apoiar a revisão pela
equipa do Hospital da Covilhã.

\section*{Turnos}

\begin{center}
\begin{tabular}{cll}
\toprule
\textbf{Código} & \textbf{Horário} & \textbf{Descrição} \\
\midrule
M  & 08:00--16:00 & Manhã \\
A  & 16:00--00:00 & Tarde normal \\
EA & 14:00--22:00 & Tarde antecipada \\
N  & 00:00--08:00 & Noite \\
\bottomrule
\end{tabular}
\end{center}

Todos os turnos contam como oito horas.

\section*{Regras obrigatórias}

\begin{enumerate}
    \item Cada técnico realiza, no máximo, um turno por dia.
    \item Técnicos em férias ou indisponíveis não podem ser escalados.
    \item Cobertura mínima de segunda a sexta-feira: 7 M, 3 A/EA e 1 N.
    \item Cobertura mínima ao sábado e domingo: 3 M, 3 A e 1 N.
    \item Em cada dia útil, exatamente uma pessoa da tarde realiza EA.
    \item Cada técnico só pode realizar os turnos para os quais está elegível.
    \item Limite máximo de 40 horas por semana ISO.
    \item Limite mensal proporcional ao contrato individual de 35 ou 40 horas,
          arredondado para turnos completos e ajustado pelas horas transitadas.
    \item A não pode ser seguido de M no dia seguinte; EA pode ser seguido de M.
    \item Depois de N, o técnico não trabalha no dia seguinte.
    \item Dois turnos N do mesmo técnico devem iniciar-se com pelo menos seis dias
          de intervalo.
    \item Cada técnico elegível recebe obrigatoriamente um número de noites entre
          a média arredondada por defeito e por excesso. Em maio de 2025: 2 ou 3 noites.
    \item Quando Angelo ou Nuno estão sem trabalhar, Lina deve realizar M, salvo
          férias de Lina. Atualmente esta regra também se aplica aos fins de semana.
    \item Angelo e Nuno alternam a prevenção diariamente: existe exatamente um P
          por dia, que pode coexistir com M e é então apresentado como M/P.
    \item Cada prevenção ao domingo ou feriado gera uma folga compensatória FC,
          colocada num dia útil sem férias nem turno regular.
\end{enumerate}

\section*{Elegibilidade individual atual}

\begin{center}
\begin{tabular}{>{\raggedright\arraybackslash}p{0.27\textwidth}
                >{\raggedright\arraybackslash}p{0.62\textwidth}}
\toprule
\textbf{Técnico} & \textbf{Turnos permitidos} \\
\midrule
Celia L & Apenas M. \\
Fernanda e Sandra & M ou A/EA; sem noites. \\
Celso & M ou A/EA, apenas em dias úteis. \\
Angelo e Nuno & Apenas M em dias úteis. \\
Restantes técnicos & M, A/EA ou N, em qualquer dia. \\
\bottomrule
\end{tabular}
\end{center}

\section*{Preferências otimizadas}

Estas regras podem ser violadas quando necessário, mas o sistema tenta:

\begin{itemize}
    \item equilibrar as horas relativamente aos dias disponíveis;
    \item evitar mais de cinco dias consecutivos de trabalho;
    \item juntar duas ou mais folgas e evitar uma folga isolada entre turnos;
    \item evitar um único turno entre dois dias de folga;
    \item igualar os dias acumulados de trabalho, compensando diferenças no mês seguinte;
    \item rodar feriados, Natal, Ano Novo e Páscoa com base no histórico;
    \item equilibrar os fins de semana relativamente às oportunidades de cada técnico;
    \item evitar fins de semana consecutivos;
    \item evitar trabalhar apenas sábado ou apenas domingo;
    \item evitar fins de semana próximos de férias;
    \item minimizar alterações quando uma escala publicada precisa de ser reparada.
\end{itemize}

\section*{Pontos prioritários para decisão}

\begin{itemize}
    \item \validar{} --- coberturas mínimas e tratamento dos feriados.
    \item \validar{} --- elegibilidade e respetiva justificação para cada técnico.
    \item \validar{} --- fórmula das horas mensais e horas contratuais individuais.
    \item \validar{} --- descanso A--M, EA--M e após N.
    \item \validar{} --- intervalo mínimo entre noites.
    \item \validar{} --- distribuição obrigatória de 2 ou 3 noites.
    \item \validar{} --- regra de Lina, Angelo e Nuno, sobretudo ao fim de semana.
    \item \validar{} --- ciclo de cinco turnos, duas folgas e tratamento das folgas adicionais.
    \item \validar{} --- histórico anterior de dias trabalhados e de feriados.
    \item \validar{} --- proteção de fins de semana antes e depois de férias.
    \item \validar{} --- prioridades entre carga horária, noites, fins de semana,
          preferências pessoais e estabilidade da escala.
\end{itemize}

\section*{Elementos ainda não modelados}

Confirmação das horas contratuais individuais; regimes de tempo parcial; qualificações por
modalidade ou função; níveis de experiência; preferências individuais; pedidos
de folga prioritários; regras de feriados; tipos distintos de ausência; horas
extraordinárias; composição mínima das equipas; e regras completas entre meses.

\section*{Registo da validação}

\small
\renewcommand{\arraystretch}{1.35}
\begin{center}
\begin{tabular}{>{\raggedright\arraybackslash}p{0.42\textwidth}cccc}
\toprule
\textbf{Área} & \textbf{Correta} & \textbf{Alterar} & \textbf{Discutir} & \textbf{N/A} \\
\midrule
Turnos e horários & $\square$ & $\square$ & $\square$ & $\square$ \\
Cobertura mínima & $\square$ & $\square$ & $\square$ & $\square$ \\
Elegibilidade individual & $\square$ & $\square$ & $\square$ & $\square$ \\
Limites de horas & $\square$ & $\square$ & $\square$ & $\square$ \\
Descanso e noites & $\square$ & $\square$ & $\square$ & $\square$ \\
Regra de Lina, Angelo e Nuno & $\square$ & $\square$ & $\square$ & $\square$ \\
Fins de semana e férias & $\square$ & $\square$ & $\square$ & $\square$ \\
Prioridade dos objetivos & $\square$ & $\square$ & $\square$ & $\square$ \\
\bottomrule
\end{tabular}
\end{center}
\normalsize

Nome e função: \hrulefill \hfill Data: \hrulefill

\vspace{0.6cm}
Observações: \hrulefill

\vspace{0.8cm}
\hrulefill

\end{document}
